\scanpage{426}
\mspage{165}
\deleted{\illegible{1}} Donc on définit un objet \deleted{de $C(c,u,v)$ \unsure{$\mid$}}
\[
(c',w')=\psi(c'',w'')\in\mathrm{Ob}\,C(c,u,v)
\]
par
\[
c^{\unsure{\prime|\prime\prime}}=c'',\qquad w'=\rho_c w'',
\]
et c'est fonctoriel en $(c'',w'')$.

Le composé $\varphi\psi$ étant un endomorphisme de la catégorie contractile $\underline{\mathrm{Ch}}(a',b')_{/\widetilde c}$, est homotope à l'identité. Il faut vérifier qu'il en est de $=$ de $\psi\varphi$, endomorphisme de $C(c,u,v)$. Ce foncteur associe à $(c',w')$, $(\widetilde{c'},\overline{w'})$, où $\widetilde{c'}$ est défini comme plus haut $\widetilde{c'}=a'_{\Delta_+}\circ c'\circ b'_{\Delta_-}$, et où $\overline{w'}$ est le composé
\diagram{\begin{tikzcd}[ampersand replacement=\&,column sep=small,row sep=small]
\widetilde{c'}\arrow[rr,"\rho_{c'}",""{name=rhoanchor,below}]\arrow[d,equals] \&\& c'\arrow[r,"w'"]\&c\\
\delta_{\tau'}^{\unsure{*|'*}}(c') \& \& {\substack{\text{\unsure{flèche}}\\\text{canonique}}} \arrow[from=rhoanchor,to=2-3,dash,bend left=12] \&
\end{tikzcd}}

On a donc
\[
(\psi\varphi)(c',w')=(\widetilde{c'},\overline{w'})\xrightarrow{\rho_{c'}}(c',w')
\]
c'est bien un morphisme dans $C(c,u,v)$ i.e. on a commutativité \unsure{du|des}
\diagram{\begin{tikzcd}[ampersand replacement=\&]
\widetilde{c'}\arrow[rr,"\rho_{c'}"]\arrow[dr,"\overline{w'}"'] \&\& c'\arrow[dl,"w'"]\&{\text{(flèche de }\mathrm{Ch}(a',b')}\\
\&c\&\&\unsure{\text{--}}
\end{tikzcd}}
% archive sheet number: 425
